% Generated draft: requires mathematical and source review.
\documentclass[11pt]{article}
\usepackage[margin=25mm]{geometry}
\usepackage{fontspec}
\setmainfont{Latin Modern Roman}
\usepackage{amsmath,amssymb,mathtools,unicode-math}
\setmathfont{Latin Modern Math}
\usepackage{enumitem,needspace,xurl,hyperref}
\setlength{\parindent}{0pt}
\setlength{\parskip}{5pt}
\setlength{\emergencystretch}{3em}
\newcommand{\N}{\mathbb{N}}
\newcommand{\Z}{\mathbb{Z}}
\newcommand{\Q}{\mathbb{Q}}
\newcommand{\R}{\mathbb{R}}
\newcommand{\C}{\mathbb{C}}
\newcommand{\F}{\mathbb{F}}
\newcommand{\A}{\mathbb{A}}
\newcommand{\PP}{\mathbb{P}}
\newcommand{\E}{\mathbb{E}}
\newcommand{\eps}{\varepsilon}
\newcommand{\dd}{\,\mathrm{d}}
\begin{document}
\section*{UnsolvedMath ID: 800145}
\textbf{Source problem code:} catalogue entry
\subsection{Definitions and mathematical statement}
Under the usual interpretation of the notation, let
\[
\psi(x)=\sum_{\substack{p\text{ prime},\ k\in\Z_{\ge1}\\p^k\le x}}\log p,
\qquad E(x)=\psi(x)-x\qquad(x\in\R,\ x\ge1),
\]
where \(\log\) denotes the natural logarithm. The supplied record asks for progress on sign changes of \(E(x)\) beyond repeated-averaging techniques. It describes an expectation that large-scale intervals of positive and negative values alternate at approximately the rate suggested by repeated averaging, while each detected transition contains many smaller-scale sign changes. It further suggests that a quantity denoted \(W^\psi(T)\), and related sign-change counting functions, should grow roughly like \(\sqrt T\), up to logarithmic factors, rather than like \(\log T\). These are supplied heuristic expectations, not asserted bounds. The record does not define \(W^\psi(T)\), the parameter \(T\), the sign-change counting convention, the scales involved, or the repeated-averaging benchmark. Consequently, this portion does not specify a precise quantified conjecture.

The same record also contains a question attributed to Alexandre Bailleul, who attributes its origin to Nilotpal Sinha. For an integer \(n\ge2\), define its largest prime factor by
\[
P^+(n)=\max\{p:p\text{ is prime and }p\mid n\}.
\]
Let the real Dirichlet character modulo \(4\) be
\[
\chi_{-4}(m)=
\begin{cases}
0,&m\equiv0,2\pmod4,\\
1,&m\equiv1\pmod4,\\
-1,&m\equiv3\pmod4
\end{cases}
\qquad(m\in\Z).
\]
For real \(x\ge2\), the well-defined part of the displayed sum is
\[
S(x)=\sum_{\substack{n\in\Z\\2\le n\le x}}\chi_{-4}\bigl(P^+(n)\bigr).
\]
Equivalently,
\[
S(x)=\#\{n\in\Z:2\le n\le x,\ P^+(n)\equiv1\pmod4\}
-\#\{n\in\Z:2\le n\le x,\ P^+(n)\equiv3\pmod4\}.
\]
The question in this passage is to determine the asymptotic behaviour of \(S(x)\) as \(x\to\infty\). The record reports numerical evidence of a bias in favour of the second set, but specifies neither a quantitative conjecture nor a precise meaning of bias. Its original notation \(\sum_{n\le x}\) leaves the lower endpoint implicit and does not define \(P^+(1)\); the displayed definition of \(S(x)\) explicitly restricts to \(n\ge2\). The record does not clarify whether this largest-prime-factor question belongs to the titled problem or was included from a separate source item.
\subsection{Short English statement}
The record asks for improved understanding of sign changes of \(\psi(x)-x\), but omits definitions needed to state its suggested growth law precisely. It also includes a possibly separate question about the asymptotic difference between counts of integers whose largest prime factor is congruent to \(1\) or \(3\) modulo \(4\). Both questions require clarification before a single precise conjecture can be extracted.
\subsection{Sources}
\begin{itemize}
\item Local UnsolvedMath registry, permanent ID \(800145\), source code \texttt{catalogue entry}, title ``Comparative Prime Number Theory --- Problem 2.'' The supplied registry labels the record open; this label is not an independently verified mathematical assessment.
\item Hamieh, Kadiri, Martin, and Ng, \emph{Comparative Prime Number Theory Problem List} (2024), supplied source item ``Problem 2 (\#14),'' \url{https://arxiv.org/abs/2407.03530}. The registry lists these authors as proposers and gives proposed year 2024. The supplied provenance records access on 2026-10-16 and extraction by \texttt{arxiv\_latex}; these are supplied provenance claims, not independently verified access or source verification.
\item The largest-prime-factor passage is attributed in the supplied text to Alexandre Bailleul, Universit\'e Paris--Saclay, and says that the question was raised by Nilotpal Sinha, with reference label [Sinha]. No bibliographic details or link for [Sinha] are supplied. The claim that the asymptotic behaviour remains unknown and the report of numerical bias are supplied source assessments, not independent literature conclusions.
\end{itemize}
\end{document}
